A continuación, se detallan las mejoras que se consideran necesarias para el futuro.

En primer lugar, se considera necesario la posibilidad de editar valoraciones ya realizadas del usuario por 
si se ha cometido algún error al introducir la puntuación o el comentario. Para ello, se podría implementar un botón 
de edición en la lista de valoraciones del usuario, que al pulsarlo,
se muestre un diálogo con el formulario de valoración ya rellenado. También, se valora que los usuarios puedan 
responde a las valoraciones de otros usuarios, para fomentar la interacción entre ellos.

En segundo lugar, se considera la posibilidad de añadir un sistema de idioma para la aplicación, para que el usuario 
pueda elegir el idioma. Para ello, se podría implementar un selector de idiomas en la barra de navegación que permita al usuario
cambiar el idioma. 

En tercer lugar, se podría implementar un sistema de edicion de la información del usuario, para que pueda cambiar su nombre de usuario, correo electrónico
o contraseña. Incluso se podrían ampliar los campos del usuario para que pueda añadir una foto de perfil o una breve biografía. 

En cuarto lugar, Angular ofrece la librería NGRX \cite{ngrx} que permite gestionar el estado de la aplicación de forma más eficiente. De tal forma que 
los filtros que se aplicaron en una pantalla se mantengan al navegar a otra pantalla y volver a la anterior, sin necesidad de 
volver a aplicar los filtros.

En quinto lugar, la implementación de un servicio caché permitiría almacenar los datos recibidos del backend y evitar 
realizar peticiones innecesarias al servidor lo que mejoraría al rendimiento de la aplicación.

Por último, se valora implementar un sistema de notificaciones para que el usuario sepa cuando ha recibido una solicitud de amistad, 
cuando ha recibido una respuesta de una valoración o para informar de la actividad de sus amigos. 